% !TEX root = ../../../main.tex
\section{Norms and angles}
\label{sec:linear-algebra-i:inner-product-norms-and-angles}

% Sources: October 22 handwritten PDF, p. 4 and top of p. 6.
\begin{definition}
  The norm induced by a real inner product is the nonnegative number
  \[
    \lVert v\rVert=\sqrt{(v,v)}.
  \]
\end{definition}

For the ordinary dot product on \(\R^n\), this is the Euclidean
magnitude. It is often convenient to work with its square,
\[
  \lVert v\rVert^2=(v,v),\qquad
  \lVert v\rVert=0\quad\Longleftrightarrow\quad v=0.
\]
% Only the norm induced by an inner product is supplied. Do not add a
% general norm-axiom definition, triangle inequality, or other norm examples.

\begin{theorem}[Cauchy--Schwarz inequality]
  \label{thm:linear-algebra-i:cauchy-schwarz}
  For any \(v,w\) in a real inner-product space,
  \[
    \lvert(v,w)\rvert\le\lVert v\rVert\lVert w\rVert.
  \]
  Equality holds if and only if the pair \(v,w\) is linearly dependent,
  including the cases in which either vector is zero.
\end{theorem}
% On October 28 Lucas authorized this standard form, correcting p. 5's
% missing absolute value and incomplete equality condition.
% The source postpones the proof. Homework 7 Problem 4.3 asks for the
% squared version, with both equality alternatives. Lucas supplied its proof
% on October 30 and authorized completion of the converse for w=sv.

% Keep the complete three-part question with its minimization setup.
\newpage
\subsection{A minimization problem}

Minimizing a squared norm leads to the Cauchy--Schwarz inequality in
\autoref{ex:hw07:minimization-and-cauchy-schwarz}. In this exercise,
single bars denote the induced norm.
% !TEX root = ../../hw07.tex
% Source: ma-gy7033fall2026homework7 (1).pdf, pp. 1--2, Problem 4 (4.1--4.3).
% Preserve the source's single-bar norm notation and its zero-vector cases.
\begin{exercise}
  \label{ex:hw07:minimization-and-cauchy-schwarz}
  (30 points) Let \(V\) be a real inner product space and let
  \(v,w\in V\) be nonzero vectors. Consider the function
  \(f:\R\to\R\) given by
  \[
    f(t)=|v-tw|^2.
  \]
  \begin{enumerate}[{4}.1.][3]
    \item (10 points) Prove that \(f\) has a unique minimum.
    \item (10 points) Find the minimum value of \(f\) and the value of
      \(t\) for which \(f(t)\) equals the minimum value.
    \item (10 points) Use your solution to the previous parts of this
      problem to prove that for any \(v,w\in V\),
      \[
        (v,w)^2\le(v,v)(w,w)
      \]
      and equality holds if and only if there exists a scalar
      \(s\in\R\) such that
      \[
        v=sw\quad\text{or}\quad w=sv.
      \]
  \end{enumerate}
  \solutionlink{hw07:minimization-and-cauchy-schwarz}
\end{exercise}


\subsection{Angles}

For nonzero \(v,w\), the angle \(\theta\in[0,\pi]\) is described by
\[
  (v,w)=\lVert v\rVert\lVert w\rVert\cos\theta,
  \qquad
  \cos\theta=\frac{(v,w)}{\lVert v\rVert\lVert w\rVert}.
\]
This relation gives a notion of angle using the inner product, including
when the vectors are functions or matrices. In the ordinary real plane,
it has the geometric meaning shown in
\autoref{fig:linear-algebra-i:inner-product-angle}.

\begin{marginfigure}
  \centering
  % Author-authorized illustration of October 22 pp. 4/6's cosine relation.
\begin{tikzpicture}[
  x=11mm,y=11mm,font=\footnotesize,line width=0.45pt,
  vector/.style={-{Latex[length=1.6mm,width=1.1mm]}}
]
  \coordinate (O) at (0,0);
  \coordinate (V) at (2.8,0);
  \coordinate (W) at (1.4,1.8);
  \draw[vector] (O) -- (V) node[right=3pt] {\(v\)};
  \draw[vector] (O) -- (W) node[above right=2pt] {\(w\)};
  \node[below=4pt] at (1.5,0) {\(\lVert v\rVert\)};
  \node[above left=3pt] at (0.8,1.03) {\(\lVert w\rVert\)};
  \draw[black!55] (0.8,0) arc[start angle=0,end angle=52.125016,radius=0.8];
  \node at (26.062508:1.08) {\(\theta\)};
  \node[below left=2pt] at (O) {\(0\)};
  \path (current bounding box.south) ++(0,-2mm);
\end{tikzpicture}

  \caption[Lengths and angle from an inner product]{A real-plane
    schematic with the ordinary dot product. The vectors have lengths
    \(\lVert v\rVert\), \(\lVert w\rVert\), and angle \(\theta\).}
  \label{fig:linear-algebra-i:inner-product-angle}
\end{marginfigure}
